\begin{frame}{Vai trò và chức năng}
\begin{table}
\centering
\caption{Bảng vai trò và chức năng trong hệ thống} % <-- Lệnh caption được thêm ở đây

% Bây giờ cả hai cột đều có thể dùng 'l' (căn trái) hoặc 'c' (căn giữa)
\begin{tabular}{|l|l|}
\hline
\textbf{Vai trò} & \textbf{Chức năng} \\
\hline

% Dòng 1: Guest
Guest &
\makecell[l]{ % [l] để căn trái các dòng bên trong ô
Giới thiệu, nghe thử giọng nói, xem gói VIP \\
Đăng ký email,
Đăng nhập \textbf{
\textcolor{blue}{G}%
\textcolor{red}{o}%
\textcolor{yellow}{o}%
\textcolor{blue}{g}%
\textcolor{green}{l}%
\textcolor{red}{e}
}\\
Xem cuộc trò chuyện được chia sẻ
} \\
\hline

% Dòng 2: User
User &
\makecell[l]{
Đăng nhập và quản lý tài khoản, cài đặt \\
Thực hiện trò chuyện, đánh dấu và chia sẻ \\
Mua và xem lịch sử gói thuê bao VIP
} \\
\hline

% Dòng 3: VIP
VIP &
\makecell[l]{
Suy luận nâng cao\\
Tài liệu cá nhân \\
Sử dụng chức năng giọng nói
} \\
\hline

% Dòng 4: Admin
Admin &
\makecell[l]{
Quản lý nhân vật,
gói VIP \\
Báo cáo đường ống dữ liệu
} \\
\hline
\end{tabular}
\end{table}

\note{

Tiếp theo em xin giới thiệu về "Vai trò và chức năng trong hệ thống"

Hệ thống gồm có 4 vai trò người dùng

Khách vãng lai +
người dùng cơ bản +
VIP +
Admin

Khách vãng lai:
Xem thông tin giới thiệu,
đăng ký email
và
đăng nhập google

Người dùng cơ bản:
Đăng nhập + quản lý thông tin,
trò chuyện + đánh dấu + chia sẻ,
mua và xem lịch sử gói thuê bao VIP

Người dùng VIP:
Sau khi người dùng mua gói thuê bao VIP thành công
sẽ có thể thực hiện các chức năng
suy luận nâng cao,
tài liệu cá nhân,
chức năng giọng nói

Quản trị viên:
Quản lý các chức năng hệ thống

nhân vật,
gói thuê bao VIP,
báo cáo đường ống dữ liệu

}
\end{frame}
